% Esquisse d'un programme (cote 119, copy 1): archive cote 119, site page 13 (PDF page 14), typescript page 11.
% typescript. Draft compared: EsquisseFr.pdf (text layer):403-478.
% AI-transcribed by Claude Opus 5.5 (claude-opus-5-5, adaptive thinking, effort high), Message Batches API, batch esq1 (msgbatch_01LNsogFyF3y6KSp9UbtJtBk), 2026-11-12; not yet checked by eye.
\scanpage{119/13}
\mspage{11}

(avec
\[ \rho_s = \sigma_2\sigma_1 \ , \quad \rho_f = \sigma_1\sigma_0 \ , \quad \sigma = \sigma_0\sigma_2 = \sigma_2\sigma_0 \ , \]
opérations de \uline{rotation élémentaire} d'un repère autour d'un sommet,
d'une face et d'une arête respectivement). Il y a un dictionnaire
parfait entre la situation topologique des cartes compactes, resp.
cartes compactes orientées, d'une part, et les ensembles finis à
groupe d'opérateurs $\underline{C}_2$ resp. $\underline{C}_2^+$ de l'autre, dictionnaire dont
l'existence était d'ailleurs plus ou moins connue, mais jamais
énoncée avec la précision nécessaire, ni développé tant soit peu.
Ce travail de fondements est fait avec le soin qu'il mérite dans
un excellent travail de DEA, fait en commun par Jean Malgoire et
Christine Voisin en 1976.

Cette réflexion prend soudain une dimension nouvelle,
avec cette remarque simple que le groupe $\underline{C}_2^+$ peut s'interpréter
comme un quotient du groupe fondamental d'une sphère orientée pri-
vée de trois points, numérotés 0, 1, 2, les opérations $\rho_s$, $\sigma$, $\rho_f$
s'interprétant comme les lacets autour de ces points, satisfaisant
la relation familière
\[ \hand{\ell}_0\hand{\ell}_1\hand{\ell}_2 = 1 \ , \]
alors que la relation supplémentaire $\sigma^2 = 1$ i.e. $\hand{\ell}_1^{\,2} = 1$ signi-%
fie qu'on s'intéresse au quotient du groupe fondamental correspon-%
dant à un indice de ramification imposé 2 au point 1, qui classi-%
fie donc les revêtements de la sphère, ramifiés au plus en les
points 0, 1, 2, avec une ramification égale à 1 ou 2 en les points
au-dessus de 1. Ainsi, les cartes orientées compactes forment une
catégorie isotopique équivalente à celle de ces revêtements, soumis
de plus à la condition supplémentaire d'être des revêtements finis.
Prenant maintenant comme sphère de référence la sphère de Riemann,
ou droite projective complexe, rigidifiée par les trois point 0, 1
et $\infty$ (ce dernier remplaçant donc 2), et se rappelant que tout revê-%
tement ramifié fini d'une courbe algébrique complexe hérite lui-même
d'une structure de courbe algébrique complexe, on aboutit à cette
constatation, qui huit ans après me paraît encore toujours aussi
extraordinaire : \uline{toute carte orientée} ``\uline{finie}'' \uline{se réalise canoniquement}
\uline{sur une courbe algébrique complexe} ! Mieux encore, comme la droite

\margin{13}
